\section{Exact triangle reduction and learned target grouping}
\label{sec:analytic-batch-method}
Some integral structures admit a direct reduction formula. We use such a formula before asking a learned model to choose the remaining native work. The new exact case is a massless one-loop triangle with two lightlike external momenta and one nonzero invariant. Its denominator powers can include a polynomial numerator. Its singular invariant matrix prevents use of the preceding inverse-matrix recurrence in the full triangle sector. Appendix~\ref{app:analytic-batch} derives the coefficient of a single off-shell bubble master from Feynman parameters and expresses it through finite products. This avoids equation generation and elimination for that block. Factorization then applies the formula independently to products of these blocks. The tests through thirteen loops concern these independent products; they do not establish a formula for a general connected thirteen-loop scattering graph.

For the coupled two-loop boxes, several requested integrals often need overlapping equations. We therefore train a model to select how Kira shares the reduction work. The five choices are the preceding application, a common reduction with either of two sparse seed regions, separate reductions for groups of equal numerator rank, and separate reductions for each target. Numerator rank is the sum of the absolute values of negative indices. Each choice uses the same target integrals and checked master basis. Exact algebra determines the coefficients, and an incomplete restricted reduction invokes the existing recovery procedure.

The model uses seventeen numerical features of the denominators and target indices. These include target count, numerator ranks, distances between numerator indices, and how propagators depend on the loop momenta. It predicts the complete cost of each choice. A cost-sensitive tree selects splits that reduce the sum of measured action costs, and bootstrap ensembles provide an additional candidate model~\cite{CART1984,Breiman1996Bagging}. Appendix~\ref{app:analytic-batch} derives the training objective and states all candidate settings. Reference coefficients enter only the independent result checks.

We collect 32 training requests and sixteen separate validation requests from planar and nonplanar massless boxes. Their indices have unit positive propagator powers and numerator ranks one through four. We select the model using validation costs, then test its full application cost against the best fixed action at a fresh validation profile. This second check includes loading, prediction, exact reduction, and verification. The model enters deployment only if it saves at least two percent, completes every validation check, and has a request-resampling interval below equal time. The reserved requests remain unused during both decisions.

During development, repeated basis discovery caused avoidable failures on some planar batches. We add a separately checked registry containing only master indices. Two native reductions with different integral orderings must reproduce each other and the development coefficients. The same-basis control receives this registry too, so it separates basis preparation from the new triangle formula and learned grouping. The unmodified preceding application supplies a further practical control. We report basis preparation and benchmark-reference preparation separately from the supplied-basis application timer.
